\chapter{Конструкторская часть}
%<пишем про алгоритмы, архитектуру ПО и тп>

\section{Описание алгоритмов}

\begin{figure}[!htbp]
	\centering
	\includesvg[width=0.7\textwidth]{aa2-rec}
	\caption{Схема рекурсивного алгоритма.}
\end{figure}

\begin{figure}[!htbp]
	\centering
	\includesvg[width=0.2\textwidth]{aa2-nonrec}
	\caption{Схема нерекурсивного алгоритма.}
\end{figure}
\clearpage

\FloatBarrier


\section{Расчёт трудоёмкости алгоритмов}

Введём модель вычислений для определения теоретической трудоёмкости описанных ранее алгоритмов.\\

\begin{table}[!htbp]
	\caption{Цены базовых операций.}
	\centering
	\begin{tabularx}{\textwidth}{|*{2}{>{\centering\arraybackslash}X|}}
		\hline
		\textbf{Операции} & \textbf{Цена} \\\hline
		{ +, - , =, ++, $--$ , <<, >> } & 1 \\\hline
		{ *, /, \%, *=, /= } & 2 \\\hline
		{ ==, >=, <=, !=, >, <, \&\&, || } & 1 \\\hline
		{ [], \{\} } & 1 \\\hline
		{ call/return } & 34 \\\hline
	\end{tabularx}
\end{table}

Цена вызова-возврата рассчитывалась по формуле \ref{cr}:
\begin{align}
	\label{cr}
	CR = 2 \cdot (p + k + r + l + f), % cr/2 = 1 + 16 + 0 + 0 + 0 = 17
\end{align}
где
\begin{itemize}[label={}]
\item $p$ - количество передаваемых параметров,
\item $k$ - количество сохранямых в стеке регистров,
\item $r$ - количество возвращаемых по адресной ссылке значений,
\item $l$ - количество локальных ячеек для обеспечения реентерабельности (повторного вызова до завершения предыдущего),
\item $f$ - количество ячеек для возвращаемого значения.
\end{itemize}

\begin{table}[!htbp]
	\centering
	\caption{Трудоёмкость алгоритмов.}
	\small
	\begin{tabularx}{\textwidth}{|*{3}{X|}}
		\hline
		\textbf{Алгоритм} & \textbf{Трудоёмкость ХС} \\\hline
		Рекурсивный & $(2N - 1)\mathrel{CR} + 3N$ \\\hline
		Нерекурсивный & $\mathrel{CR} + 2 + 3N$ \\\hline 
	\end{tabularx}
\end{table}

\FloatBarrier

\section*{Вывод}
%<что сделали в конструкторке и получили в результате, кратко>

В данном разделе мной были построены схемы алгоритмов и вычислена их трудоёмкость с помощью введённой модели вычислений.

\clearpage
